% JSAI 2027 (41st Annual Conference) general-session paper, English. Félix's target: 4 pages
% of body text, references after. Derived from ../main.tex, which is left untouched.
% Reframed 2026-10-17 after the mock RSS review panel: every failure runs to the time limit,
% so a calibrated timer catches all failures and the endpoint is warning lead over the timer.
% Numbers and conventions match ../rss2027_stage1 (analysis/lead_time.py, extra.py; an alarm
% at the final inference counts). JSAI-only extras: analysis/jsai_extras.py.
% The conformal proposition and the Stage-2 plan stay RSS-only.
%
% TEMPLATE NOTE: the official JSAI English LaTeX package (jsaiac_tex_en-*.zip) could not
% be downloaded from this environment, so this preamble reproduces the published rules
% (A4, two columns, centred title block, abstract <= 150 words centred above the columns,
% contact address at the bottom of the first column). Before submission, move the body
% (everything between BEGIN BODY and END BODY) into the official template's main file and
% fill its title/author/affiliation/address macros with the values below.
%
% Deadlines (JSAI 2027, from the venue plan): abstract 2027-02-03, manuscript
% 2027-02-17 14:00 JST, PDF only, <= 3 MB. See the dual-submission note in README.md.
\documentclass[10pt,a4paper,twocolumn]{article}

\usepackage[a4paper,top=20mm,bottom=22mm,left=17mm,right=17mm,columnsep=7mm]{geometry}
\usepackage[T1]{fontenc}
\usepackage{mathptmx}
\usepackage{amsmath}
\usepackage{graphicx}
\usepackage{booktabs}
\usepackage{array}
\usepackage{xcolor}
\usepackage{tikz}
\usetikzlibrary{arrows.meta,positioning}
\usepackage{pgfplots}
\pgfplotsset{compat=1.17}
\usepackage[numbers,sort&compress]{natbib}
\usepackage[hidelinks]{hyperref}
\usepackage{titlesec}
\usepackage{enumitem}

\titleformat{\section}{\normalsize\bfseries}{\thesection.}{0.5em}{}
\titleformat{\subsection}{\small\bfseries}{\thesubsection}{0.5em}{}
\titlespacing*{\section}{0pt}{1.0ex plus .3ex}{0.5ex}
\titlespacing*{\subsection}{0pt}{0.8ex plus .2ex}{0.3ex}
\setlength{\bibsep}{1pt}
\renewcommand{\bibfont}{\footnotesize}
\pagestyle{empty}

\newcommand{\pending}[1]{\textcolor{red}{[PENDING: #1]}}
\newcommand{\todo}[1]{\textcolor{orange}{[TODO: #1]}}

% ---- Title block values (copy into the official template) ----
\newcommand{\papertitle}{Time Is the Baseline to Beat:\\
Evaluating VLA Failure Detectors Against Timers}
\newcommand{\paperauthors}{Felix Doublet\textsuperscript{*1}\quad \todo{co-authors}}
\newcommand{\paperaffil}{\textsuperscript{*1}\todo{Institute of Science Tokyo, Shinoda Lab: confirm official affiliation}}
\newcommand{\paperaddress}{Contact: Felix Doublet, \todo{affiliation, postal address, e-mail}}

\begin{document}

\twocolumn[{%
\begin{center}
{\Large\bfseries \papertitle\par}
\vspace{3mm}
{\normalsize \paperauthors\par}
\vspace{1mm}
{\small \paperaffil\par}
\vspace{3mm}
\begin{minipage}{0.86\textwidth}
\small
% Abstract: <= 150 words (JSAI rule). Current count: about 145.
Online failure detectors for vision-language-action (VLA) robot policies are usually compared
by recall at a fixed false-positive rate (FPR). On RoboCasa365, a failed rollout is one that
reaches its time limit, so this comparison misses the obvious baseline: a timer. We evaluate
a learned detector against timers on 760 prospective rollouts of Xiaomi-Robotics-1 on 38
tasks. Calibrated at 5\% FPR, a single timer at 90\% of the limit flags all 380 failures, so
only warning lead time can separate detectors. Fusing the SAFE failure score with this timer
raises the mean lead from 0.090 to 0.207 of the limit (paired gain 0.117, 95\% CI
0.060--0.178). But a per-task timer, set at the 95th percentile of each task's successful
durations plus a calibrated margin, reaches 0.195 at 4.5\% FPR without reading the policy,
and SAFE adds only 0.007 to it (CI $-0.003$ to 0.019). Failure detectors should be evaluated
by their lead over a per-task timer.
\end{minipage}
\vspace{5mm}
\end{center}}]

{\renewcommand{\thefootnote}{}\footnotetext{\footnotesize \paperaddress}}

% =========================== BEGIN BODY ===========================
\section{Introduction}

Generalist vision-language-action (VLA) policies still fail often on household tasks, so a
deployed robot needs an online signal that the current attempt is going wrong. Recent
detectors read the policy's hidden state and learn to predict eventual
failure~\cite{gu2025safe}, or score how unusual that state is relative to successful
rollouts~\cite{xu2025faildetect,romer2025fiper}. They are compared with each other, usually
by ROC-AUC or by recall at a fixed false-positive rate (FPR), but rarely with what the robot
knows for free: how long the attempt has been running.

That comparison matters because of how failure is defined. In RoboCasa365, as in most
simulated manipulation benchmarks, an episode ends at first success or at a time limit. In
our prospective test cohort, every one of the 380 failed rollouts ran exactly to its task's
limit. ``Detecting'' such a failure therefore means predicting that the limit will be
reached, and a timer that alarms at 90\% of the limit already does so for all 380 failures
at 4.7\% FPR. Completed-rollout duration alone reaches ROC-AUC 0.80 on held-out rollouts,
against 0.53 for the raw SAFE score. SAFE removes this shortcut offline by truncating
rollouts to equal length~\cite{gu2025safe}, but a deployed detector sees rollouts of unequal
length, and a timer is a legitimate deployable detector.

Against the timer, a learned detector can add only one thing: an earlier alarm at the same
false-alarm rate, which leaves more time to stop, reset or ask for help. This paper is an
evaluation, not a new detector: we test one learned detector, SAFE~\cite{gu2025safe},
against a single timer and a per-task timer. Detectors are developed on one cohort, their thresholds are frozen on a separate
calibration cohort at 5\% FPR, and they are evaluated once on a disjoint prospective cohort.
Our contributions are:
\begin{enumerate}[leftmargin=*,itemsep=1pt,topsep=2pt]
  \item \textbf{The timer baseline.} A calibrated timer flags every failure within the FPR
  budget, while a survival-style time-risk curve misses 37\%, and we show why.
  \item \textbf{Lead time as the endpoint.} We measure the mean warning lead, the fraction
  of the time limit left when the first alarm fires (zero for a missed failure).
  \item \textbf{How much the policy adds.} OR-fusing the SAFE score with a single timer
  more than doubles its mean lead, and noise does not; but a per-task timer that never reads
  the policy matches it, and SAFE adds only 0.007 of the limit on top.
\end{enumerate}

\section{Related Work}

\textbf{Failure detectors for VLA policies.} SAFE trains a small network on a VLA's
last-layer features to predict eventual failure and thresholds it with a functional
conformal band~\cite{gu2025safe}; we reuse its implementation. Follow-ups change the
training data, labels or features~\cite{rajaprakash2026safecast,park2026hideseek,ma2026failurespot,yang2026triinfo,zhu2026vlascope}.
FAIL-Detect~\cite{xu2025faildetect} and FIPER~\cite{romer2025fiper} score the policy
against successful rollouts only; FIPER alarms when two scores \emph{both} cross, whereas we
alarm when \emph{either} does. Action-chunk monitors~\cite{agia2024sentinel,seligmann2026vlafail,huang2026actprobe}
compare overlapping chunks, which our policy does not produce. CheckVLA~\cite{liu2026checkvla},
the closest work on our benchmark, verifies each chunk with a world model, and
Foresight~\cite{wang2026foresight} uses world-model latents; vision-language monitors judge
progress from frames~\cite{liu2023reflect,sinha2024realtime,zhou2025codeasmonitor,duan2025aha}.
None of these works compares against elapsed time. Duration is a shortcut
feature~\cite{geirhos2020shortcut}, but its causal form is legitimate evidence online.

\textbf{Time and intervention.} Our time risk is a survival-style
estimate~\cite{kaplan1958,cox1972}, and lead time at a fixed false-alarm budget is the
accuracy--earliness trade-off of early classification~\cite{xing2012early}.
RESCUE~\cite{ahmedt2026rescue} shows that good detection scores do not guarantee lower
intervention cost, so the value of an earlier alarm must eventually be shown in closed loop.

\section{Problem Setting}

\textbf{Time.} We measure time in four units (Table~\ref{tab:time}). The simulator advances
in \emph{control steps} at 20\,Hz, RoboCasa's default control rate, which our collection
code does not change. The policy is queried once per 16-step action chunk and executes the
whole chunk, so one \emph{policy inference} covers 16 control steps, or 0.8\,s. Inferences
are our clock: $t = 1, 2, \dots$ counts them, and a detector can change its decision only
at an inference. The official time limit of task $k$ is $T_k$ control steps; in
inferences it is $H_k = \lceil T_k/16 \rceil$, between 29 and 272 for our tasks (22.5 to
217.5\,s). Videos keep every second control step and play at 20 frames per second, twice
real time; we use them only for figures. Every fraction of time in this paper, such as
$t/H_k$, the timer fraction $c$ and the lead $\ell$, is a fraction of $H_k$.

\begin{table}[t]
\centering
\caption{Units of time, with the failed \emph{OpenCabinet} rollout of
Fig.~\ref{fig:qualitative} as an example.}
\label{tab:time}
\footnotesize
\setlength{\tabcolsep}{3pt}
\begin{tabular}{@{}>{\raggedright\arraybackslash}p{0.25\columnwidth}>{\raggedright\arraybackslash}p{0.5\columnwidth}>{\raggedright\arraybackslash}p{0.17\columnwidth}@{}}
\toprule
Unit & Definition & Example \\
\midrule
Control step & one simulator action, 20\,Hz (50\,ms) & 1{,}050 \\
Policy inference $t$ & one 16-step action chunk (0.8\,s) & 66 \\
Time limit $T_k$, $H_k$ & official limit, control steps; $H_k = \lceil T_k/16\rceil$ in inferences & 1{,}050; 66 \\
Video frame & every 2nd control step, shown at 20\,fps & 526 \\
Elapsed fraction & $t/H_k$ (SAFE alarm) & 30/66 \\
Warning lead $\ell$ & $(H_k - T)/H_k$ at alarm $T$ & 36/66 = 0.55 \\
\bottomrule
\end{tabular}
\end{table}

\textbf{Monitoring.} A detector watches one rollout of a known task $k$ and decides at each
inference $t$ whether to alarm, using only the policy features so far, $t$, the task and
$H_k$; it may not use the rollout's final length or outcome, and it never changes the
policy's actions. The first threshold crossing is the alarm $T$ (in inferences). Failure is the positive class. A successful rollout with any alarm before it ends is a false
positive, and FPR is the fraction of successful rollouts with a false alarm. An alarm at any
inference up to and including the last one counts, with $\ell = 0$ at the last one. For a failed rollout, which ends at
$H_k$, the \emph{warning lead} is $\ell = (H_k - T)/H_k$ if an alarm fires and $\ell = 0$
otherwise.
We report the mean lead $\bar\ell$ over failures, which counts a miss as no warning and so
combines recall and earliness in one number. Unlike recall, it does not saturate when a
timer catches every failure, and unlike a median over detected failures, it cannot be
improved by giving up hard failures: a detector that drops a failure loses that failure's
whole lead. A timer at fraction $c$ has $\bar\ell \approx 1 - c$ by construction, up to
rounding to whole inferences, so any gain over it must come from alarms before $c\,H_k$. We also report recall, FPR, the median lead in
inferences over detected failures and the fraction of failures alarmed by $0.75\,H_k$.
Detectors are compared at operating points fixed on a calibration cohort under
FPR $\le \alpha = 0.05$. Confidence intervals come from a paired bootstrap that resamples
tasks, then rollouts within each task (2,000 replicates). Only the SAFE + timer versus timer
comparison was singled out in advance as the main comparison; all other intervals are
descriptive.

\section{Detectors}

\begin{figure}[t]
\centering
\resizebox{\columnwidth}{!}{%
\begin{tikzpicture}[font=\scriptsize,
  box/.style={draw, rounded corners=2pt, align=center, minimum height=0.85cm, inner sep=3pt},
  arr/.style={-{Latex[length=1.6mm]}, semithick}]
  \node[box, fill=gray!8] (obs) at (0,0) {Observation\\+ instruction};
  \node[box, fill=gray!8] (vla) at (2.2,0) {VLA policy};
  \node[box, fill=blue!8] (score) at (4.9,0) {Policy score $s_t$\\(SAFE MLP)};
  \node[box, fill=gray!8] (task) at (2.2,-1.35) {Task $k$, limit $H_k$,\\inference $t$};
  \node[box, fill=gray!8] (time) at (4.9,-1.35) {Timer\\$t / H_k$};
  \node[box, fill=orange!12] (or) at (7.4,-0.675) {$s_t \ge \tau_s$\\\textbf{or}\\$t/H_k \ge c$};
  \node[box, fill=red!10] (alarm) at (9.1,-0.675) {Alarm};
  \draw[arr] (obs) -- (vla);
  \draw[arr] (vla) -- node[above] {$h_t$} (score);
  \draw[arr] (task) -- (time);
  \draw[arr] (score.east) -- (or.north west);
  \draw[arr] (time.east) -- (or.south west);
  \draw[arr] (or) -- (alarm);
  \node[align=center, font=\scriptsize, text=gray!70!black] at (7.4,-1.8)
    {$(\tau_s, c)$ chosen jointly on a\\calibration cohort at FPR $\le 0.05$};
\end{tikzpicture}}
\caption{SAFE + timer. At each policy inference, the policy score and the elapsed fraction
of the time limit are compared with two thresholds; the first crossing of either raises the
alarm. The timer guarantees an alarm on every failure; the policy score can only move it
earlier.}
\label{fig:method}
\end{figure}

\textbf{Timers.} The global timer alarms once $t \ge c\,H_k$ for a fraction $c \in (0, 1]$
shared by all tasks; it uses only the task's time limit. The per-task timer uses
$c_k = \min(1, \rho_k + \delta)$, where $\rho_k$ is the 95th percentile of task $k$'s
development success durations as a fraction of $H_k$. The shared margin $\delta$ is set on the
calibration successes by a split-conformal quantile rule~\cite{angelopoulos2021gentle}, which
bounds the FPR on a new successful rollout by 0.05. It needs successful rollouts of each task,
but no failures and no policy features.

\textbf{Time risk.} The survival-style alternative estimates
\begin{equation}
  r_t = \Pr(\text{failure} \mid \text{task } k,\ \text{still running at } t)
\end{equation}
from development rollouts. With $n_t$ of them still running at $t$ and $f_t$ of those
failing, it fits an increasing isotonic regression in $t/H_k$ to $(f_t+0.5)/(n_t+1)$,
weighted by $n_t$. Because every failure runs to $H_k$, $r_t$ is in effect a task-specific
timer whose threshold is shared across tasks.

\textbf{Policy score.} At each inference we take the final-layer action-token hidden state
of the policy's flow-matching action head (1,024 dimensions). The SAFE
MLP~\cite{gu2025safe} maps it to a failure score $u^{(j)}_t$ for each of three training seeds
$j$. With per-task location $m^{(j)}_k$ and scale $\sigma^{(j)}_k$ estimated on development
rollouts, the policy score is
\begin{equation}
  s_t = \frac{1}{3}\sum_{j=1}^{3} \frac{\max_{t' \le t} u^{(j)}_{t'} - m^{(j)}_k}{\sigma^{(j)}_k}.
\end{equation}
The running maximum makes $s_t$ monotone and does not change when a threshold is first
crossed. We use the MLP head because it ranked failures best after task normalization on
held-out development data (ROC-AUC 0.69 against 0.65 for the LSTM head). We also fit four success-only scorers on the same hidden state
(Mahalanobis~\cite{lee2018mahalanobis}, cosine $k$NN~\cite{sun2022knn}, PCA-KMeans and
RND~\cite{burda2019rnd}), using successful development rollouts only.

\textbf{OR-fusion.} A fused detector alarms at the first inference where either signal
crosses its threshold, for example
\begin{equation}
  T = \min\{\, t : s_t \ge \tau_s \ \text{or}\ t \ge c\,H_k \,\}
\end{equation}
for SAFE + timer (Fig.~\ref{fig:method}). The pair of thresholds is searched jointly over the
per-rollout maxima of the calibration cohort; among pairs with calibration FPR $\le 0.05$,
we keep the one with the largest mean lead, or the highest recall for the original SAFE +
time-risk design. This selection is empirical: choosing the best of many pairs on 114
calibration successes is optimistic and carries no finite-sample FPR guarantee. SAFE +
per-task timer instead moves both thresholds along one path ($\delta = \lambda$,
$\tau_s = u + s\lambda$, with $u$ and $s$ chosen on calibration failures) and sets $\lambda$ by
the same conformal rule.

\textbf{Null fusions.} To test whether a fusion gains only because failed rollouts run
longer and give any score more chances to cross, we fuse the timer with the time risk (a
second time signal) and with Gaussian noise (the running maximum of i.i.d.\ draws, 20 draws),
and calibrate them exactly like SAFE + timer.

\section{Experimental Setup}

\textbf{Policy and tasks.} We use the released Xiaomi-Robotics-1 RoboCasa365
checkpoint~\cite{xiaomi2026robotics}. It predicts 16-step action chunks with a five-step
flow-matching action head and executes each chunk in full. We take the 39 official
RoboCasa365 tasks~\cite{nasiriany2024robocasa,robocasa365} whose published success count is
between 5 and 45 of 50, and drop one with too few failures, leaving 38 tasks (10 atomic, 28
composite). Failed rollouts last 29 to 272 inferences (median 127, about 100\,s).

\begin{table}[t]
\centering
\caption{Cohorts. S/F = successes/failures; $q$ = rollouts per class per task. No two
cohorts share an environment seed.}
\label{tab:cohorts}
\footnotesize
\begin{tabular}{@{}llr@{}}
\toprule
Cohort & Used to & Primary (S/F, $q$) \\
\midrule
Development & fit SAFE, scorers, time risk & 1,824 rollouts \\
Calibration & set thresholds & 228 (114/114, 3) \\
Prospective test & evaluate once & 760 (380/380, 10) \\
Confirmatory & preregistered test & 760 (380/380, 10) \\
\bottomrule
\end{tabular}
\end{table}

\textbf{Cohorts and protocol.} Four disjoint cohorts each serve one purpose
(Table~\ref{tab:cohorts}). Each evaluation cohort is collected from fresh environment seeds
until every task has a fixed quota of successes and failures, and the primary cohort takes
the first rollouts of each class per task in seed order. SAFE, the normalizations, the
time-risk curves and an original staged detector were frozen before the prospective test
was collected. The timer and all fused detectors were added afterwards and evaluated once
on the opened cohort, with thresholds from calibration only; each analysis first reproduces
the frozen results and checks the input checksums. The results below are therefore
exploratory, and the confirmatory cohort has not been opened.

% Qualitative example shared with the RSS draft, read from its folder so both papers stay in
% sync. Frames stay grey placeholders until frames/f1.png..f4.png exist (see its README).
% For submission, copy that folder here and change \qualdir.
\def\qualdir{../rss2027_stage1/figures/qualitative}
\input{\qualdir/fig_qualitative}

\section{Results}

\begin{table}[t]
\centering
\caption{Prospective test (380 successes, 380 failures), thresholds from calibration at
FPR $\le 0.05$. $\bar\ell$ = mean warning lead as a fraction of the time limit (miss = 0);
Med = median lead in inferences over detected failures; @.75 = fraction of failures alarmed
by $0.75\,H_k$. $^\ast$Mean over 20 draws ($\bar\ell$ range 0.088--0.134). $^\dagger$95th
percentile of development success durations, conformal margin on the 114 calibration
successes.}
\label{tab:main}
\footnotesize
\setlength{\tabcolsep}{3.5pt}
\begin{tabular}{@{}lccccc@{}}
\toprule
Detector & Recall & FPR & $\bar\ell$ & Med & @.75 \\
\midrule
Timer ($c = 0.903$) & \textbf{1.000} & 0.047 & 0.090 & 11.5 & 0.000 \\
Time risk & 0.632 & 0.047 & 0.120 & 16.5 & 0.211 \\
Time risk + timer & \textbf{1.000} & 0.050 & 0.125 & 11.0 & 0.105 \\
Noise + timer$^\ast$ & \textbf{1.000} & 0.066 & 0.101 & -- & -- \\
\midrule
SAFE only & 0.566 & 0.039 & 0.190 & 23.0 & 0.421 \\
SAFE + time risk & 0.955 & 0.058 & 0.184 & 12.0 & 0.311 \\
SAFE + timer & \textbf{1.000} & 0.058 & \textbf{0.207} & 10.0 & 0.421 \\
\midrule
Per-task timer$^\dagger$ & \textbf{1.000} & 0.045 & 0.195 & 13.5 & 0.342 \\
SAFE + per-task timer$^\dagger$ & \textbf{1.000} & 0.042 & 0.202 & 14.0 & 0.389 \\
\bottomrule
\end{tabular}
\end{table}

\textbf{The timer is the baseline to beat.} With $c = 0.903$ chosen on calibration
(FPR 0.044 there), the timer alarms on all 380 test failures at 0.047 test FPR
(Table~\ref{tab:main}). Its lead is fixed by construction: 0.090 of the limit, a mean of
11.7 inferences or about 9\,s. In an earlier version of this analysis, we measured fused
detectors against the time-risk curve and found recall gains from 0.632 to 0.90--0.96 for
SAFE and the four success-only scorers. Against the timer, those recall gains vanish: no
learned or fused detector has higher recall than 1.000, and none matches the timer's FPR at
full recall.

\textbf{Why the time risk misses.} All 140 failures the time risk misses come from the 14
tasks with at most 14 failed development rollouts. With $f$ development failures, the curve
can rise only to $(f+0.5)/(f+1)$, which stays below the shared threshold
$\tau_r = 0.972$ unless $f \ge 17$. In the 24 tasks with 17 or more development failures, the
time risk catches every failure; in the other 14, it catches none. Its misses are an
artefact of the estimator's smoothing and of a threshold shared across tasks, not a
limitation of elapsed time.

\begin{figure}[t]
\centering
\begin{tikzpicture}
\begin{axis}[width=\columnwidth, height=5.0cm, xmin=0.3, xmax=1, ymin=0, ymax=1.02,
  xlabel={elapsed fraction of the time limit $t/H_k$},
  ylabel={failures alarmed (of 380)},
  label style={font=\scriptsize}, tick label style={font=\scriptsize},
  legend style={font=\scriptsize, at={(0.02,0.98)}, anchor=north west, draw=none, fill=none},
  legend cell align=left, grid=major, grid style={gray!20}]
\addplot[gray, dashed, thick] coordinates {(0.30,0.000) (0.90,0.000) (0.92,0.842) (0.94,1.000) (1.00,1.000)};
\addlegendentry{Timer}
\addplot[black, dotted, thick] coordinates {(0.30,0.000) (0.60,0.000) (0.62,0.026) (0.64,0.026) (0.66,0.053) (0.68,0.079) (0.70,0.079) (0.72,0.079) (0.74,0.105) (0.76,0.105) (0.78,0.158) (0.80,0.158) (0.82,0.184) (0.84,0.211) (0.86,0.237) (0.88,0.237) (0.90,0.237) (0.92,0.842) (0.94,1.000) (1.00,1.000)};
\addlegendentry{Time risk + timer}
\addplot[orange!80!black] coordinates {(0.30,0.000) (0.40,0.000) (0.42,0.026) (0.44,0.026) (0.46,0.026) (0.48,0.053) (0.50,0.053) (0.52,0.053) (0.54,0.053) (0.56,0.063) (0.58,0.066) (0.60,0.074) (0.62,0.100) (0.64,0.100) (0.66,0.132) (0.68,0.161) (0.70,0.192) (0.72,0.229) (0.74,0.305) (0.76,0.318) (0.78,0.387) (0.80,0.413) (0.82,0.471) (0.84,0.513) (0.86,0.579) (0.88,0.597) (0.90,0.613) (0.92,0.684) (0.94,0.750) (0.96,0.821) (0.98,0.895) (1.00,0.955)};
\addlegendentry{SAFE + time risk}
\addplot[blue!75!black, thick] coordinates {(0.30,0.000) (0.32,0.003) (0.34,0.003) (0.36,0.011) (0.38,0.021) (0.40,0.024) (0.42,0.024) (0.44,0.024) (0.46,0.034) (0.48,0.045) (0.50,0.079) (0.52,0.103) (0.54,0.126) (0.56,0.158) (0.58,0.187) (0.60,0.221) (0.62,0.247) (0.64,0.308) (0.66,0.342) (0.68,0.355) (0.70,0.379) (0.72,0.389) (0.74,0.400) (0.76,0.424) (0.78,0.437) (0.80,0.458) (0.82,0.466) (0.84,0.474) (0.86,0.487) (0.88,0.495) (0.90,0.500) (0.92,0.500) (0.94,0.500) (0.96,0.563) (0.98,1.000) (1.00,1.000)};
\addlegendentry{SAFE + timer}
\end{axis}
\end{tikzpicture}
\caption{Cumulative fraction of the 380 prospective failures alarmed by each elapsed fraction
of the time limit, at thresholds frozen on calibration. No plotted detector alarms before
$0.3\,H_k$. The timer catches everything in one step just after $0.9\,H_k$; SAFE + timer
catches half of the failures earlier and the rest at its late timer fallback ($c = 0.96$).}
\label{fig:lead}
\end{figure}

\textbf{SAFE adds lead over a single timer.} SAFE + timer keeps recall at 1.000 and raises the
mean lead to 0.207 of the limit, a paired gain of 0.117 over the timer (95\% CI
0.060--0.178), or 20.2 against 11.7 inferences (about 16\,s against 9\,s). The SAFE
threshold fires first on 194 failures (51\%), and 42\% of failures are flagged by
$0.75\,H_k$, against none for the timer (Fig.~\ref{fig:lead}). Figure~\ref{fig:qualitative} shows one such rollout. Calibration chose a late
timer fallback ($c = 0.96$) and spent the false-alarm budget on the SAFE threshold, so the
median lead over detected failures (10 inferences) is below the timer's: the gain is
concentrated in the failures SAFE flags early, and the remaining failures are flagged later
than by the timer alone. The original SAFE + time-risk design also alarms before the timer
on 233 failures, but it misses 17.

\textbf{Sensitivity.} SAFE + timer's test FPR of 0.058 exceeds the budget (22 of 380 false
alarms, Wilson 95\% interval 0.039--0.086). Against a timer tuned on the test cohort itself to the same 0.058 FPR
($c = 0.86$, $\bar\ell = 0.136$), which favours the timer, the gain is 0.071. The result does
not depend on how an alarm at the final inference is counted: if it does not count,
calibration picks $c = 0.894$, and the timer reaches $\bar\ell = 0.104$ against 0.225 for
SAFE + timer (gain 0.121, 95\% CI 0.062--0.184).

\textbf{Where the lead comes from.} SAFE + timer warns earlier than the timer in 19 of
the 38 tasks and slightly later in the other 19 (by about 0.06, the cost of its fallback at
$c = 0.96$). The gain is largest in the tasks with few development failures, exactly those
where the time risk never fires (Spearman $\rho = -0.84$ between per-task gain and
development failure count). These are high-success tasks such as opening a cabinet.

\textbf{A per-task timer removes the gain.} The per-task timer reaches $\bar\ell = 0.195$ at
0.045 FPR without reading the policy. SAFE + timer leads it by an insignificant 0.012 (CI
$-0.044$ to $0.063$), and by $-0.003$ against a per-task timer tuned on test to the same FPR.
SAFE + per-task timer reaches 0.202, a gain of 0.007 (CI $-0.003$ to $0.019$), far below the
0.05 of the limit we would call useful. The two effects cancel across tasks: where the time
risk never fires, SAFE + timer leads the per-task timer (0.405 against 0.221); elsewhere the
per-task timer leads (0.180 against 0.092). In Fig.~\ref{fig:qualitative}, the per-task timer
alarms at $t = 38$, eight inferences after SAFE.

\textbf{Sensitivity to the percentile.} We fixed the 95th percentile after the test cohort had
been opened, so we report others. With the longest development success, the per-task timer
reaches only 0.125 and SAFE adds 0.057 (CI $-0.003$ to $0.119$); with the 75th percentile, 0.132
and 0.064 (CI 0.029--0.099); with the 90th, 0.169 and 0.028 (CI $-0.012$ to $0.067$). SAFE's
added lead therefore shrinks as the time baseline improves, and is near zero against the
best one we found.

\textbf{Time alone and noise do not reproduce it.} Fusing the time risk with the timer, which
uses every time signal we have, reaches $\bar\ell = 0.125$. Noise fused with the timer under
the same search reaches 0.101 at a higher FPR (0.066). Both fall well short of 0.207. SAFE
also separates outcomes among rollouts of equal length: among rollouts still running at 10,
25, 50 and 75\% of the limit, its within-task ROC-AUC is 0.55, 0.60, 0.59 and 0.58, against
0.51, 0.45, 0.47 and 0.41 for noise and exactly 0.5 for any time signal, which is constant
within a task at a fixed step. The hidden state thus carries some outcome information
beyond time, but at 5\% FPR a per-task timer captures nearly all the lead it adds.

\begin{table}[t]
\centering
\caption{Recall / FPR / $\bar\ell$ by task type on the prospective test (atomic: 10 tasks,
100 rollouts per class; composite: 28 tasks, 280 per class). Descriptive only: thresholds
were set on the pooled calibration cohort, so per-type FPR is not controlled.}
\label{tab:strata}
\footnotesize
\setlength{\tabcolsep}{3pt}
\begin{tabular}{@{}lcc@{}}
\toprule
Detector & Atomic & Composite \\
\midrule
Timer & 1.000 / 0.020 / 0.079 & 1.000 / 0.057 / 0.094 \\
Time risk + timer & 1.000 / 0.020 / 0.086 & 1.000 / 0.061 / 0.138 \\
SAFE + time risk & 0.980 / 0.030 / 0.217 & 0.946 / 0.068 / 0.172 \\
SAFE + timer & 1.000 / 0.090 / 0.364 & 1.000 / 0.046 / 0.151 \\
\bottomrule
\end{tabular}
\end{table}

\textbf{Atomic and composite tasks.} The lead gain is largest on atomic tasks, where SAFE +
timer reaches $\bar\ell = 0.364$ against 0.079 for the timer, but at 0.090 FPR
(Table~\ref{tab:strata}). On composite tasks, it reaches 0.151 against 0.094 at 0.046 FPR, and
the time risk + timer comes close (0.138). A plausible reason is that a terminal failure
label fits one atomic attempt, but marks the correct early work of a long composite task as
failing. A first test did not support this: on five composite tasks, SAFE trained on
segmented subtask labels ranked held-out stages no better than chance (ROC-AUC 0.49).

\textbf{Warnings remain late.} No detector alarms on more than 8\% of failures by half of the
limit. The lead that SAFE adds lies mostly in the second half of the rollout.

\pending{confirmatory test on 760 fresh rollouts (base seed 1300007), after amending its
preregistration so that the per-task timer (95th percentile) is the primary comparator and
$\Delta\bar\ell$ the primary endpoint.}

\section{Discussion and Limitations}

On benchmarks where failure means reaching a time limit, a calibrated timer already detects
every failure, so a learned VLA failure detector should be credited only for warning earlier
than a timer at the same false-alarm rate. Against a single timer, OR-fusing the SAFE score
more than doubles the mean lead, and noise does not. Against a per-task timer built from
successful durations, SAFE adds 0.007 of the limit. On our data, what looked like policy
evidence is mostly time calibration.

The finding also concerns benchmark design: when every failure is a timeout, recall at a
fixed FPR rewards any score that grows with elapsed time. Papers on VLA failure detection
should report a per-task timer and the lead each detector gains over it, and benchmarks
should label failures that become unrecoverable before the time limit.

\textbf{Limitations.} (i)~The timer and all fused detectors were defined after the
prospective cohort had been opened, so these results are exploratory until the confirmatory
test. (ii)~The two-threshold search is optimistic: fused detectors reach 0.050--0.066 test
FPR against 0.044 on calibration, with only 114 calibration successes; a conformal
calibration of the OR rule and a larger calibration cohort would address this.
(iii)~The gain is a dozen inferences on average, mostly late in the rollout, and whether an
earlier alarm raises task success through reset and retry is untested. (iv)~The simulator
ends failed episodes only at the time limit, so failures that become unrecoverable earlier
(drops, collisions) are not labelled as such. (v)~SAFE needs white-box access to the
policy's hidden state and per-task development data that includes failures; the timer
needs only the time limit. (vi)~All results use one policy in simulation, and the task and
its time limit are assumed known. (vii)~The per-task timer's percentile was chosen after the
cohort was opened; with other percentiles, SAFE adds 0.03--0.06 of the limit.

\section{Conclusion}

Elapsed time is the baseline a VLA failure detector must beat. When failures run to the time
limit, a calibrated timer catches every failure within a 5\% FPR budget, so learned detectors
should be judged by lead time. On 760 prospective rollouts across 38 RoboCasa365 tasks, SAFE
fused with a single timer raised the mean warning lead from 0.090 to 0.207 of the limit, but
a per-task timer built from successful durations reached 0.195 without reading the policy,
and SAFE added 0.007 to it. Next steps are a preregistered test against the per-task timer, a
second policy, and failures that end before the time limit.

% ============================ END BODY ============================
\clearpage

\bibliographystyle{unsrtnat}
\bibliography{refs}

\end{document}
